\documentclass{article}
\usepackage{amsmath,amssymb,amsthm}
\usepackage{graphicx,booktabs,xspace}
\usepackage[numbers]{natbib}
\newtheorem{thm}{Theorem}
\newtheorem{lem}[thm]{Lemma}
\newcommand{\eg}{e.g.\xspace}

\title{Detecting Clock Drift in Low-Cost Sensor Networks}
\author{A. N. Author}
\date{}

\begin{document}
\maketitle

\begin{abstract}

Cheap sensors keep time with cheap oscillators, and their clocks drift by
several seconds per day. When readings from many sensors are merged, this
drift shifts events in time and blurs the picture of how pollution moves
through a city. We propose a method that detects and corrects drift after
the fact, using only the readings themselves. The idea is that sensors close
to each other see the same events; the time lag between their versions of an
event estimates their relative drift. On a network of 24 sensors, the method
recovers the drift of each unit to within 0.4 seconds per day, which allows
to align the readings without any extra hardware.

\end{abstract}

\section{Introduction}

Every reading a sensor sends comes with a timestamp, and most analyses take
that timestamp at face value. For a single sensor this is harmless. For a
network it is not: an event that passes two sensors one minute apart is
recorded with a lag of one minute plus the difference of their clock errors.
If the clocks drift, the lag drifts with them, and estimates of wind speed
or of the direction in which a plume travels become unreliable.

The usual remedy is to synchronize the clocks over the network, \eg with NTP
or a GPS receiver \citep{elson2002fine}. Both are often unavailable on
low-cost units. Many deployments use radios that only send and cannot
receive, and a GPS receiver adds cost and drains the battery. As a
consequence, much of the data that low-cost networks collect carries
timestamps that nobody can trust to better than a few minutes.

We take a different route and estimate the drift afterwards, from the data.
Concretely, we use the fact that neighbouring sensors observe the same
events: a delivery truck, a bonfire, the morning traffic. The lag between
two copies of the same event has a real part, the travel time of the air,
and a part that is due to the clocks. Over many events the real part
averages to a constant, while the clock part grows linearly in time.
Separating the two gives the relative drift of every pair of sensors; a
least-squares fit turns the pairs into one estimate per sensor. The method
needs no extra hardware and works with less sensors than one might expect,
since three neighbours per sensor are enough in our experiments.

\section{Related Work}

Clock synchronization in sensor networks is a mature topic. Protocols such
as reference broadcast \citep{elson2002fine} and flooding-based schemes
\citep{maroti2004flooding} reach microsecond accuracy, but they all need
nodes that can receive messages. Post-hoc methods are less common. Some
align records by matching them to an external signal, such as the mains
frequency in audio recordings; others use known events, such as a scheduled
test alarm. Our setting has neither, and we rely on events that nobody
scheduled.

Lag estimation by cross-correlation is standard in acoustics and in
seismology, where it is used to locate sources. We borrow the tool, not the
goal: we do not locate events, and we treat the travel time of the air as a
nuisance rather than as the quantity of interest.

\section{Method}

In what follows, sensor $i$ reports the time $\tau_i(t) = t + a_i + b_i t$,
where $a_i$ is an offset and $b_i$ the drift rate. We want to estimate $b_i$
for every sensor. The offsets matter less, since most analyses only compare
readings within a day.

For two neighbouring sensors $i$ and $j$, we cut the readings into windows
of one hour and find, in each window, the lag that maximizes the
cross-correlation of the two series. Let $\ell_{ij}(w)$ be that lag in
window $w$. Recall that the lag contains the travel time of the air, which
we treat as noise with a constant mean. We then fit

\begin{equation}
  \ell_{ij}(w) = c_{ij} + (b_i - b_j)\, t_w + \varepsilon_{ij}(w)
  \label{eq:lag}
\end{equation}

by least squares, which gives an estimate of the difference $b_i - b_j$ for
every pair. Put differently, each pair of neighbours contributes one linear
equation in the unknown drift rates.

The pairwise equations determine the drift rates only up to a common
constant. To fix it, we need one sensor with a trusted clock, or we assume
that the drift rates average to zero across the network. We use the second
option unless a reference station is available; the error it introduces is
the mean drift of the network, which is small when the units come from
different production batches.

\begin{thm}\label{thm:var}
  If each window lag has variance at most $\sigma^2$, the least-squares estimate of $b_i - b_j$ from
  $n$ windows spanning a time $T$ has variance at most $12 \sigma^2 / (n T^2)$.
\end{thm}

The bound explains why long recordings help so much: the variance falls with
the square of the recording length. A week of data is enough for most pairs
in our network. The proof is the standard argument for the slope of a
least-squares line and is given in Appendix~\ref{app:var}.

\section{Experiments}

We installed 24 sensors in a residential area, each within 300 meters of at
least three others. Six of them received a GPS module that recorded the true
time; these serve as ground truth and are hidden from the method. Despite of
the simple setup, the network had the usual problems of field work: two
units lost power for several days, and one was moved by a resident.

Table~\ref{tab:drift} compares the estimated drift of the six GPS units with
their true drift. The mean absolute error is 0.4 seconds per day, against
true drift rates between 2 and 9 seconds per day. The error is largest for
the unit that was moved, since its neighbours changed halfway through the
recording\footnote{We kept this unit in the evaluation on purpose; dropping
it lowers the mean error to 0.3 seconds per day.}.

\begin{table}[t]
  \centering
  \caption{Estimated and true drift of the six units with GPS (seconds per day).}
  \label{tab:drift}
  \begin{tabular}{lcccccc}
    \toprule
    Unit & G1 & G2 & G3 & G4 & G5 & G6 \\
    \midrule
    True & 2.1 & 4.8 & 9.0 & 3.3 & 6.2 & 5.5 \\
    Estimated & 2.4 & 4.6 & 8.7 & 3.1 & 7.0 & 5.6 \\
    \bottomrule
  \end{tabular}
\end{table}

We also checked how the accuracy depends on the number of neighbours. With
two neighbours the error doubles; with four or more it barely changes. We
stress that the neighbours do not have to be close in a strict sense. Any
sensor that sees the same events at a stable lag will do.

The window length is the one setting that needs care. Windows shorter than
half an hour contain too few events, and the lag estimates become noisy;
windows longer than three hours average over changes in the wind. One hour
worked well in our data. To be safe, we recommend checking the stability of
the lag estimates on a few pairs before running the full fit.

\section{Discussion}

The method assumes that the drift rate is constant over the recording. Real
oscillators also drift with temperature, so in a climate with strong daily
cycles the estimate is an average over the cycle. A natural extension fits a
separate rate for each temperature band; we leave this to future work.

A second assumption is that events are shared. Sensors in very different
places, such as one inside a building and one on a roof, see different
events, and their lag estimates are noise. In practice we drop pairs whose
correlation peak is weak. This leaves a sparser graph, but a more reliable
one.

\section{Conclusion}

Clock drift is a quiet source of error in low-cost sensor networks. We
showed that it can be estimated from the readings themselves, without extra
hardware, to within a fraction of a second per day. The method is simple
enough to run on a laptop for a network of a few hundred sensors.

\bibliographystyle{plainnat}
\bibliography{refs}

\appendix
\section{Proof of Theorem~\ref{thm:var}}\label{app:var}

The estimate of $b_i - b_j$ is the slope of a least-squares line through $n$
points whose abscissas are spread evenly over a time $T$. Its variance is
$\sigma^2$ divided by the sum of squared deviations of the abscissas, which
is at least $n T^2 / 12$. This gives the bound.

\end{document}
